\documentclass{article}
\usepackage[T1]{fontenc}
\usepackage{lmodern}
\usepackage{graphicx}
\usepackage{amsmath,amssymb}
\usepackage[a4paper,margin=2cm]{geometry}
\usepackage[absolute,overlay]{textpos}
\setlength{\TPHorizModule}{1mm}
\setlength{\TPVertModule}{1mm}
\usepackage[indonesian]{babel}
\usepackage{hyperref}
\setlength{\emergencystretch}{2em}
% unit: Halaman 11

\begin{document}

% === Halaman 11 (llm) ===

\section*{Parameter umum Diffie-Hellman}

\begin{itemize}
    \item Misalkan dua entitas yang berkomunikasi adalah Alice dan Bob.
    \item Mula-mula Alice dan Bob menyepakati sebuah bilangan prima $p$, dan sebuah bilangan bulat $g$, sedemikian sehingga $g < p$ dan $g$ adalah akar primitif dari $p$.
    \item Nilai $p$ dan $g$ tidak perlu rahasia. Bahkan, Alice dan Bob dapat membicarakannya melalui saluran publik yang tidak aman sekalipun.
\end{itemize}

\vspace{10mm}

\begin{center}
Rinaldi M/II4020 Kriptografi
\end{center}

\end{document}
